% ---------------------------------------------------------------------------------------------
% Appendix: three-action extension (PACA-RPS v1). \input this file after the last appendix.
% Needs: fig_rps_temp1_results.pdf (upload to figures/) and table_rps_temp1.tex (same folder as this file).
% Cross-references used: fig:results (main Figure 2), app:opponent (opponent appendix),
%                        app:temperature (the new prompt-by-temperature appendix; rename to match).
% All numbers are from rps/results/ (prespecified analysis, tag rps-v1), requested temperature 1 unless noted.
% ---------------------------------------------------------------------------------------------
\section{Extension to a three-action game}\label{app:rps}

\paragraph{Design.} To test whether the assay extends beyond a two-action game, the same seven models played standard rock--paper--scissors against a three-action version of Proteus. The prompts, opponent, schedule, and analysis were fixed and version-tagged before collection. Each model completed 16 sessions of 200 rounds (112 sessions, 22,400 rounds), run as adjacent minimal/informed pairs that shared an opponent seed, with four pairs at each of two requested temperatures (0 and 1) and prompt order counterbalanced within each temperature. The minimal prompt asked the model to choose the object that defeats the one on which the blue circle lands and stated which object defeats which. The informed prompt added that the opponent analyzes previous choices, prescribed choosing each action with probability $1/3$ independently across rounds, and warned that alternation, repetition, and cycling can be exploited. History display, response format, and parsing followed the matching-pennies procedure. A win scored $+1$, a loss $-1$, and a tie $0$, so independent uniform play earns 0 in expectation against any opponent that commits before seeing the current choice. The opponent kept the nine channels of Appendix~\ref{app:opponent}. Each channel tested whether any action followed the current context at a rate different from $1/3$ (exact binomial tests at $p<.05$, Bonferroni-corrected over the three actions). Proteus then played the best response to the channel with the largest expected gain $g$ with probability $g$, and otherwise chose uniformly; with two actions this rule reduces to the binary opponent. It is an adaptation of the algorithm of \citet{Lee2004}, not a published algorithm. In 1,000 simulated sessions against an independent uniform player, mean payoff was $-.002$ (95\% CI $[-.006, +.002]$) and at least one channel was active on 9.9\% of rounds, compared with 10.7\% for the binary opponent. All 112 sessions completed with no substituted choices. Three attempts by Gemini 3 Flash Preview ended in an unparseable response and were rerun from round 1 with the same seed, as prespecified.

\paragraph{Proteus exploited history dependence that marginal frequencies missed.} At requested temperature 1 (Figure~\ref{fig:rps}, Table~\ref{tab:rps}), mean payoff was negative in all 14 model-by-prompt cells and in 53 of 56 sessions. Averaged over prompts, models earned $-.253$ per round (95\% CI $[-.368, -.137]$, $t(6) = -5.36$, $p = .002$), losing 50.1\% of rounds and winning 24.8\%. An opponent choosing uniformly at random holds any player to an expected payoff of 0, so these losses measure what the adaptive opponent extracted from regularities in the models' play. As in matching pennies, the largest losses often came with balanced action frequencies. Under the informed prompt, Sonnet and Opus kept every action within .02 of $1/3$, closer than independent play typically produces, but repeated their previous action on only 1\% and 2\% of rounds, and lost $.49$ and $.33$ per round. Under the minimal prompt, Gemini 3.1 Flash Lite was similarly balanced (marginal bias $.024$) but on 68\% of rounds chose the action that its previous action defeats, and it had the lowest payoff in the study ($-.63$). Other losses reflected marginal bias: GPT-5.4 chose SCISSORS on 70\% of informed rounds and lost $.38$ per round. In an exploratory regression across the 14 cells, marginal bias alone accounted for almost none of the variation in payoff ($R^2 = .01$), and adding excess mutual information between consecutive actions raised this to $R^2 = .43$. Results at requested temperature 0 were similar (all 14 cell means negative; mean payoff $-.297$).

\paragraph{The effect of informed prompting was model-specific.} The prespecified primary test asked whether informed prompting lowered payoff at requested temperature 1, as it lowered win rate in matching pennies. It did not on average: the mean informed-minus-minimal difference was $+.027$ (95\% CI $[-.251, +.305]$, $t(6) = 0.24$, $p = .82$), with a decline in three of seven models, and the estimate at requested temperature 0 was $-.003$ (95\% CI $[-.346, +.339]$). The average conceals large effects of opposite sign. Informed prompting lowered payoff for Sonnet ($-.445$) and Opus ($-.190$), which played close to the benchmark under the minimal prompt and almost never repeated an action under the informed one, the three-action counterpart of their over-alternation in matching pennies. It raised payoff for Gemini 3.1 Flash Lite ($+.526$), Gemini 3 Flash Preview ($+.144$), and GPT-5.4 ($+.105$), which were heavily exploited under the minimal prompt and exchanged sequential dependence for marginal bias under the informed one. Unlike in matching pennies, then, minimal prompting did not bring these models close to equilibrium play. Four limits apply. The minimal prompt states which object defeats which, so it discloses the game, although not the opponent or the equilibrium. Calibration showed that this opponent gains little from mild under-repetition (a repeat rate of $.20$ cost a simulated player about $.02$ per round). Each cell has four sessions, compared with seven in the main study. Finally, the temperature setting was not applied for Opus 4.8, GPT-5.4, and GPT-5.4-mini (Appendix~\ref{app:temperature}), so for those models the two temperature conditions are repeated samples at the provider default. We conclude that the assay and its central diagnostic point carry over to a three-action game, whereas the direction of the prompt effect depends on the model there; the finding that prescribing the equilibrium strategy impairs its execution is therefore established for matching pennies and, in rock--paper--scissors, for Sonnet and Opus only.

\begin{figure}[htbp]
\centering
\includegraphics{figures/fig_rps_temp1_results.pdf}
\caption{Rock--paper--scissors against Proteus at requested temperature 1. (A) Payoff per round, with models in the order of Figure~\ref{fig:results}. (B) Payoff against marginal bias, the total-variation distance between the action frequencies and $(1/3, 1/3, 1/3)$. (C) Payoff against excess mutual information between consecutive actions. Gray symbols indicate minimal prompting and orange symbols informed prompting; large symbols are means of four sessions and faint symbols individual sessions. The dashed line marks the expected payoff of independent uniform play (0). Shaded bands give the central 95\% range of a four-session mean under independent uniform play, estimated from 40,000 simulated sessions; they are reference ranges, not evidence of equivalence. Labels name the informed-prompt point of each model, and one minimal-prompt point is labeled in gray.}
\label{fig:rps}
\end{figure}

\begin{table}[htbp]
\centering
\caption{Rock--paper--scissors against Proteus: payoff and choice structure by model and prompt at requested temperature 1.}
\label{tab:rps}
\footnotesize
\setlength{\tabcolsep}{3.6pt}
\begin{tabular}{@{}lrrrrrrrrrr@{}}
\toprule
 & \multicolumn{2}{c}{Payoff per round} & \multicolumn{2}{c}{Informed $-$ minimal} & \multicolumn{2}{c}{Marginal bias} & \multicolumn{2}{c}{Excess MI} & \multicolumn{2}{c}{Repeat rate} \\
\cmidrule(lr){2-3}\cmidrule(lr){4-5}\cmidrule(lr){6-7}\cmidrule(lr){8-9}\cmidrule(l){10-11}
Model & Min. & Inf. & $T=1$ & $T=0$ & Min. & Inf. & Min. & Inf. & Min. & Inf. \\
\midrule
Sonnet 4.6 & $-.049$ & $-.494$ & $-.445$ & $-.473$ & $.028$ & $.011$ & $.043$ & $.562$ & $.45$ & $.01$ \\
Haiku 4.5 & $-.159$ & $-.181$ & $-.023$ & $+.060$ & $.041$ & $.067$ & $.023$ & $.230$ & $.30$ & $.11$ \\
Gemini 3.1 FL & $-.626$ & $-.100$ & $+.526$ & $+.571$ & $.024$ & $.128$ & $.356$ & $.146$ & $.17$ & $.18$ \\
GPT-5.4 & $-.480$ & $-.375$ & $+.105$ & $+.191$ & $.039$ & $.365$ & $.158$ & $.040$ & $.23$ & $.48$ \\
GPT-5.4-mini & $-.083$ & $-.010$ & $+.073$ & $+.070$ & $.090$ & $.101$ & $.069$ & $.016$ & $.21$ & $.28$ \\
Opus 4.8 & $-.136$ & $-.326$ & $-.190$ & $-.478$ & $.038$ & $.015$ & $.088$ & $.479$ & $.50$ & $.02$ \\
Gemini 3 FP & $-.330$ & $-.186$ & $+.144$ & $+.035$ & $.039$ & $.262$ & $.070$ & $.045$ & $.25$ & $.34$ \\
\midrule
Mean & $-.266$ & $-.239$ & $+.027$ & $-.003$ & $.043$ & $.136$ & $.115$ & $.217$ & $.30$ & $.20$ \\
\bottomrule
\end{tabular}
\par\smallskip\begin{minipage}{\linewidth}\footnotesize Cells are means of four 200-round sessions; models are in the order of Figure~\ref{fig:results}. Payoff is $+1$ for a win, $-1$ for a loss, and $0$ for a tie, so independent uniform play earns 0 in expectation. Differences are means of four within-pair differences at each requested temperature $T$, and the bottom row averages over the seven models; the 95\% confidence intervals of the mean differences, from $t(6)$, are [$-.251$, $+.305$] at $T=1$ and [$-.346$, $+.339$] at $T=0$. Marginal bias is the total-variation distance between the action frequencies and $(1/3, 1/3, 1/3)$, which with three actions equals the largest deviation of any frequency from $1/3$. Excess MI is the mutual information between consecutive actions, in bits, minus its mean under independent uniform play. Under independent uniform play a four-session mean falls in $[.021, .062]$ for marginal bias and $[-.008, +.012]$ for excess MI with 95\% probability, and the repeat rate is $1/3$.\end{minipage}
\end{table}

